% !TEX program = xelatex
% Build (from this directory):
%   xelatex -interaction=nonstopmode shaogefenhao-ai-product-termio.tex
%   (run twice if TOC/refs need it)
\documentclass[aspectratio=169,10pt,UTF8]{ctexbeamer}

\usetheme{metropolis}
\usecolortheme{default}

\setbeamertemplate{navigation symbols}{}
\setbeamertemplate{footline}[frame number]

\usepackage{booktabs}
\usepackage{graphicx}
\usepackage{hyperref}
\usepackage{listings}
\usepackage{tikz}
\usetikzlibrary{positioning}

\lstset{
  basicstyle=\ttfamily\footnotesize,
  breaklines=true,
  columns=fullflexible,
  frame=single,
  backgroundcolor=\color{black!3},
}

\title{AI 编程范式下产品开发全过程}
\subtitle{termio.sh · Terminal-first Agentic Development Environment}
\author{袁继伟}
\date{少个分号 · 40--60 min}
\institute{\url{https://termio.sh}}

\begin{document}

\begin{frame}[plain]
  \titlepage
\end{frame}

\begin{frame}{大纲}
  \begin{enumerate}
    \item \textbf{序章} — Demo 演示
    \item \textbf{缘起} — 问题驱动：为什么做 Terminal-first ADE？
    \item \textbf{产品设计} — 用户画像、核心场景、MVP
    \item \textbf{技术实现} — 架构总览 + 三个关键难点
    \item \textbf{Agentic Engineering} — skills · workflow
    \item \textbf{迭代优化} — 从可用到好用
    \item \textbf{未来展望}
    \item \textbf{总结与启示}
  \end{enumerate}
\end{frame}

% ---------- Demo ----------
\begin{frame}{序章：Demo 演示}
  先看产品，再听故事。

  \vspace{0.8em}
  \begin{tabular}{@{}llp{5.5cm}@{}}
    \toprule
    段 & 看什么 & 证明什么 \\
    \midrule
    1 & Claude / Codex / OpenCode 并排 & 多 agent 编排，不是又一个 IDE \\
    2 & 侧栏状态点 + 菜单栏 tray & idle / working / needsAttention / done \\
    3 & iPhone 扫码续会话 & 本地 PTY host + companion \\
    \bottomrule
  \end{tabular}
\end{frame}

\begin{frame}{一句话定位}
  \begin{quote}
    \Large Stop babysitting terminals.\\[0.4em]
    \normalsize 在真实原生终端里，编排你已经在用的 coding agent。
  \end{quote}
  \vspace{0.6em}
  \begin{itemize}
    \item 不是 Cursor 插件 · 不是 Electron 套壳 · 不是云端托管 agent
    \item \textbf{是}：macOS 原生 · libghostty · 本地 PTY · 可选 iOS 远程
  \end{itemize}
\end{frame}

% ---------- 缘起 ----------
\begin{frame}{缘起 —— 问题驱动}
  Agent 时代，最慢的一环变了：
  \begin{enumerate}
    \item \textbf{窗口 babysitting} — 十几个 terminal，不知道谁在跑
    \item \textbf{状态不可见} — 思考中还是卡权限？
    \item \textbf{工具碎片} — Claude Code / Codex / OpenCode / Pi \ldots
    \item \textbf{桌面 $\leftrightarrow$ 手机断层}
  \end{enumerate}
  \vspace{0.5em}
  \centering 问题不是「不会写代码」，是 \textbf{编排与注意力}。
\end{frame}

\begin{frame}{为什么是 Terminal-first？}
  \begin{itemize}
    \item Coding agent 的「家」天然在终端（CLI 已是主流）
    \item TUI 就是 agent 的 UI —— 再包一层 Web IDE 是重复劳动
    \item 终端 = 真实 PTY = 真实权限 / shell / 工具链
  \end{itemize}
  \vspace{1em}
  \centering
  \textbf{IDE 帮人写代码；ADE 帮人管 agent。}
\end{frame}

\begin{frame}{关键产品判断}
  \begin{columns}[T]
    \begin{column}{0.48\textwidth}
      \textbf{做}
      \begin{itemize}
        \item 原生 macOS（Swift）
        \item 拥有 PTY（libghostty）
        \item 多 agent + 状态 tray
        \item 本地优先、无账号
        \item 配置驱动适配 agent
      \end{itemize}
    \end{column}
    \begin{column}{0.48\textwidth}
      \textbf{不做}
      \begin{itemize}
        \item Electron / 纯 Web 壳
        \item 只 tail 日志、不拥有运行时
        \item 做成完整 IDE
        \item 强制云同步 / 账号门禁
        \item 每加 agent 改一堆 Swift
      \end{itemize}
    \end{column}
  \end{columns}
\end{frame}

\begin{frame}{赛道位置（2026-07）}
  同桌面：Unpeel · cmux · Conductor · herdr · Warp Oz \ldots

  \vspace{0.6em}
  \textbf{termio 押注的差异}
  \begin{enumerate}
    \item 原生 + libghostty + 零配置状态 + 菜单栏 ambient
    \item first-class worktree / 会话编排
    \item iOS companion（监督面，不是第二套 TUI）
    \item 小而意见化 —— 不和 cmux 拼 breadth
  \end{enumerate}
\end{frame}

% ---------- 产品 ----------
\begin{frame}{用户画像（只打一个核心）}
  \textbf{重度 CLI agent 用户}
  \begin{itemize}
    \item 日常用 Claude Code / Codex / OpenCode
    \item 多 session 并行，靠自己记「谁在干活」
    \item 在乎：本地、隐私、原生手感、键盘流
    \item 不在乎：再学一个云 IDE、再开一个账号
  \end{itemize}
\end{frame}

\begin{frame}{核心场景 + MVP}
  \textbf{三个场景}
  \begin{enumerate}
    \item 启动与并行（project $\to$ multi-session）
    \item 盯状态（侧栏 + tray）
    \item 离开工位仍在线（iOS companion）
  \end{enumerate}
  \vspace{0.5em}
  \textbf{MVP 必须有}：侧栏 · 真 PTY · agent 预设 · 状态可见 · 无账号\\
  \textbf{故意后置}：第二个 Cursor · never-die host · 过早平台化
\end{frame}

\begin{frame}{关键产品决策（真实做过的）}
  \begin{itemize}
    \item 2026-06-26：Zed fork \textbf{推倒重来} $\to$ Swift + libghostty
    \item Agent = \textbf{配置文件}，不是枚举（加 Grok 曾改 6 个文件）
    \item 状态三层：OSC/铃 $\to$ hooks $\to$ title 规则
    \item Git 面板做到 TextKit 级 diff（对照 Desktop / JetBrains）
    \item iOS：有终端，但不把 TUI 当唯一监督面
  \end{itemize}
\end{frame}

% ---------- 技术 ----------
\begin{frame}[fragile]{端到端架构}
\begin{lstlisting}[basicstyle=\ttfamily\tiny]
macOS termio (Swift + AppKit/SwiftUI)
  Sidebar · TermioStore · MenuBar · Git · Editor
       |
  libghostty (Metal) <-> PTY <-> shell / agent CLI
  HookListener <- termio agent report (unix socket)
  AgentCatalog  <- Resources/agents/*.json + ~/.termio/
  CompanionServer (WebSocket)
       |  TermioShared.WireProtocol
       v
  iOS TermioMobile          Landing / Docs (Next.js · R2)
\end{lstlisting}
\end{frame}

\begin{frame}{难点 1：终端设计}
  \textbf{目标}：原生、快、会话可切换、agent TUI 全屏可用。
  \begin{itemize}
    \item libghostty（Metal），非 WebView
    \item SurfaceCache：切换 session 重挂 surface，不重 spawn
    \item PTY 生命周期 $\times$ focus / resize
  \end{itemize}
  \vspace{0.4em}
  \textbf{近三周真实坑}：窄 grid banner 冻住 · resize 不 reflow · focus 丢失
\end{frame}

\begin{frame}{难点 2：Agent 适配}
  \textbf{抽象} \texttt{AgentDefinition}：id · command · resume · icon · hook/title rules\\
  同一加载路径：bundled JSON + 用户 manifest

  \vspace{0.5em}
  \textbf{状态}：\texttt{idle $\to$ working $\to$ needsAttention | done}\\
  权威：hooks · OSC 铃 · title 规则（仲裁）

  \vspace{0.5em}
  \textbf{public contract}\\[0.3em]
  \texttt{termio agent report <working|attention|done|idle>}

  \vspace{0.4em}
  channel-scoped socket —— dev 不能盖写 release hooks
\end{frame}

\begin{frame}{难点 3：iOS 移动端}
  \begin{itemize}
    \item Mac = \textbf{PTY host}（唯一真相源）
    \item iOS = companion client
    \item \texttt{WireProtocol}：auth · attach · start · resize · files · trace
  \end{itemize}
  \vspace{0.5em}
  近三周：Chats/Projects · FAB · Markdown 同步 · renderer 稳定性 · TestFlight

  \vspace{0.5em}
  \centering
  \textbf{手机需要终端能力，但不该把 agent TUI 当唯一 UI。}
\end{frame}

\begin{frame}{三个难点的共性}
  \begin{tabular}{@{}lll@{}}
    \toprule
    难点 & 核心矛盾 & 解法 \\
    \midrule
    终端 & 真 PTY vs 易嵌入 & 拥有 surface + 修上游 \\
    Agent & CLI 碎片 vs 统一体验 & Manifest + hook + 多层状态 \\
    iOS & 远程可见 vs 本地真相 & Mac host + 薄 wire \\
    \bottomrule
  \end{tabular}
  \vspace{0.8em}
  原则：本地优先 · 零配置兜底 · 配置驱动 · 先可用再精确
\end{frame}

% ---------- Agentic Engineering ----------
\begin{frame}{Agentic Engineering：回路}
\begin{center}
\begin{tikzpicture}[
  box/.style={draw, rounded corners, align=center, minimum width=2.6cm, minimum height=0.7cm, font=\small},
  arr/.style={->, thick}
]
  \node[box] (a) {目标 / 取舍};
  \node[box, below=0.35cm of a] (b) {docs design/bug/rfc};
  \node[box, below=0.35cm of b] (c) {skills / prompts};
  \node[box, below=0.35cm of c] (d) {worktree + 并行 agent};
  \node[box, below=0.35cm of d] (e) {commit · rebuild · 验证};
  \node[box, below=0.35cm of e] (f) {用户反馈};
  \draw[arr] (a) -- (b); \draw[arr] (b) -- (c);
  \draw[arr] (c) -- (d); \draw[arr] (d) -- (e);
  \draw[arr] (e) -- (f); \draw[arr] (f.west) -- ++(-0.6,0) |- (b.west);
\end{tikzpicture}
\end{center}
\end{frame}

\begin{frame}{Skills（仓库里真实存在）}
  \begin{tabular}{@{}ll@{}}
    \toprule
    Skill & 干什么 \\
    \midrule
    \texttt{macos-rebuild-dev} & 杀进程 $\to$ 构建 $\to$ 拉起 dev app \\
    \texttt{ios-rebuild-dev} & 装真机 / 模拟器 \\
    \texttt{bump-version} & 打 tag 发版 \\
    \texttt{conventional-commit} & 规范提交 \\
    \texttt{doc} & 设计文档 + wiki 索引 \\
    \texttt{app-screenshot-debug} & 点 UI + 截窗 \\
    \bottomrule
  \end{tabular}
  \vspace{0.6em}
  高频易错步骤 $\to$ skill；一次性探索 $\to$ 普通对话
\end{frame}

\begin{frame}{可带走的清单}
  \textbf{值得固化}
  \begin{itemize}
    \item 构建与重装路径（Mac / iOS）
    \item 发版、commit 规范
    \item 设计决策与事故 handoff
    \item 带验收标准的任务
  \end{itemize}
  \textbf{不要过度自动化}
  \begin{itemize}
    \item 还在摇摆的产品取舍 · UI 审美判断 · 证书类一次性操作
  \end{itemize}
  \centering AI 加速回路；\textbf{判断仍在人。}
\end{frame}

% ---------- 迭代 ----------
\begin{frame}{迭代：近三周真实主题}
  \small
  \begin{tabular}{@{}llp{4.2cm}@{}}
    \toprule
    主题 & 信号 & 演进 \\
    \midrule
    Agent 状态 & 转完还在转 / 黄点绿点 & hooks · title · done \\
    \texttt{/clear} & trace 旧数据 & conversation rotation \\
    Git & 不像 Desktop/WebStorm & TextKit diff \\
    Agent 扩展 & 加 agent 改多文件 & \texttt{agents/*.json} \\
    iOS & FAB · TestFlight & companion parity \\
    终端细节 & 换行 · focus · reflow & 布局时序 + 上游 \\
    \bottomrule
  \end{tabular}
\end{frame}

\begin{frame}{从可用到好用}
  \begin{enumerate}
    \item 状态可信是 ADE 的生命线
    \item 原生手感是付费意愿的前置
    \item 移动端是监督，不是缩小版桌面
    \item 删掉错误抽象也是进度（如自研 sandbox）
    \item 文档是异步团队协议（一人 + 多 agent 也需要）
  \end{enumerate}
\end{frame}

% ---------- 展望 / 总结 ----------
\begin{frame}{未来展望（可讨论）}
  \begin{enumerate}
    \item 瓶颈：写代码 $\to$ 编排与确认
    \item ADE 与 IDE 长期并存
    \item 拥有运行时的人赢监督
    \item 本地 + 移动监督是专业用户默认形态
    \item Agent 编排 agent（Sessions MCP 等）
  \end{enumerate}
\end{frame}

\begin{frame}{总结与启示}
  \begin{enumerate}
    \item 从痛点出发，场景压到 3 个以内 —— MVP 是裁
    \item 先选对形态，再堆功能
    \item 为 AI 协作修路：doc · skills · worktree · 验收
    \item 优先打磨信任相关体验（状态 / 性能 / 像素）
    \item 保持删的勇气
  \end{enumerate}
\end{frame}

\begin{frame}{收束}
\begin{center}
\begin{tikzpicture}[
  n/.style={draw, rounded corners, align=center, font=\small, minimum height=0.65cm, minimum width=3.2cm},
  a/.style={->, thick}
]
  \node[n] (1) {痛点 $\to$ 定位 $\to$ MVP};
  \node[n, below=0.3cm of 1] (2) {架构三个难点};
  \node[n, below=0.3cm of 2] (3) {skills / docs / 并行 agent};
  \node[n, below=0.3cm of 3] (4) {用户反馈 $\to$ 再裁 / 再磨};
  \node[n, below=0.3cm of 4] (5) {用户愿意每天打开};
  \foreach \x/\y in {1/2,2/3,3/4,4/5} {\draw[a] (\x) -- (\y);}
\end{tikzpicture}

\vspace{0.8em}
\textbf{产品 $\times$ 工程 $\times$ AI 回路 = 从想法到喜爱}
\end{center}
\end{frame}

\begin{frame}[plain]
  \centering
  \vspace{1.5cm}
  {\LARGE 谢谢}\\[1em]
  {\large Termio — the terminal home for your AI coding agents}\\[0.6em]
  \url{https://termio.sh}\\[1.5em]
  {\Large Q \& A}
\end{frame}

\end{document}
